% 358_Harmonik_Algebra_En_ch.tex
% Johann Pascher, 8 September 2026
% Harmonics and algebraic structure: why the same mathematics — and what it yields
\providecommand{\BM}{\textbf{[B]}}
\providecommand{\KM}{\textbf{[K]}}
\providecommand{\SM}{\textbf{[S]}}
\providecommand{\XM}{\textbf{[X]}}
\providecommand{\QM}{\textbf{[Q]}}
\providecommand{\EM}{\textbf{[E]}}

\phantomsection
\addcontentsline{toc}{chapter}{%
  Harmonics and Algebraic Structure --- Why the Same Mathematics}

\section*{Status markers}
\begin{center}
\begin{tabular}{@{}lp{10.5cm}@{}}
\textbf{[K]} & \emph{Confirmed} --- numerically verified. \\[3pt]
\textbf{[B]} & \emph{Proved} --- algebraically established. \\[3pt]
\textbf{[S]} & \emph{Speculative} --- open. \\[3pt]
\textbf{[E]} & \emph{Established} --- textbook mathematics, not FFGFT-specific. \\[3pt]
\textbf{[Q]} & \emph{Source} --- external primary value. \\
\end{tabular}
\end{center}

\section*{Abstract}

The corpus describes harmonic structures with the algebra of finite
fields in several places: commas and fractal correction (Doc.~060),
overtones as torus spectrum (Doc.~159), the entry rule of harmonic primes
and chords as partitions (Doc.~342), zeta hierarchy and resolution floor
(Doc.~343), consonance as coupling geometry (Doc.~328). This document does
not repeat those results. It answers two questions left open there:
\emph{why} the same mathematics describes both subjects --- and
\emph{what that yields}.

\textbf{Results:}
\begin{enumerate}
\item Theorem~A \BM: In a finite field every ``circle of fifths'' closes
  exactly; a comma is impossible. The musical comma and the fractal
  correction $K_\text{frak}$ are therefore not properties of the Galois
  structure but of the embedding into the continuum. This explains why
  $K_\text{frak}$ becomes visible in the corpus only at the SI transition
  (Doc.~333).
\item Theorem~B \BM: The two exceptions of the 5-limit reading in Doc.~060
  (strange with prime 13, top with prime 7) are not exceptions in the
  Galois reading: 13 and 7 are the harmonic primes of depth $k=3$ and
  $k=6$ (Doc.~343 Theorem~E). All nine Yukawa coefficients lie in the
  Galois grid $\{2,3,5,7,11,13\}$; the prime 11 ($k=5$) does not occur.
\item Observation~C \SM: Three twelves coincide --- twelve-tone
  temperament as the quotient of the Tonnetz by two commas,
  $|(\mathbb{Z}/13)^*|=12$ as the carrier group of the four fermion
  sectors, and the Farey denominator height $Q_c\approx11$--$12.5$ at
  comma resolution. The mechanisms differ; the agreement in order is an
  observation, not a theorem.
\item Theorem~D \BM: The answer to the why-question: $\tilde{T}\cdot m=1$
  turns mass ratios into period ratios; period ratios are characters of
  cyclic groups; the $Z_3$ reduction turns these into arithmetic in
  GF$(3^k)$. Harmonics and Galois classification are not analogous but
  two representations of the same object.
\end{enumerate}

Verification script: \texttt{pruef\_358\_harmonik\_algebra.py} (22/22).

% ============================================================
\section{Starting point: what the corpus already says}
% ============================================================

The following table assigns the existing results and is only cited below.

\begin{center}
\begin{tabular}{@{}lll@{}}
\toprule
Harmonics & Algebra / FFGFT & Doc. \\
\midrule
Overtone series $f, 2f, 3f,\dots$ & Laplace spectrum on $S^1$, $T^d$ & 159 \\
Pythagorean comma $+1.364\,\%$ & $K_\text{frak}=1-100\xi$, $-1.333\,\%$ & 060, 133 \\
5-limit ratios & Yukawa coefficients $r_i$ & 006, 060 \\
Circle of fifths & torus topology, winding numbers & 060, 189 \\
Temperament & renormalisation / $K_\text{frak}$ & 060 \\
Consonance $\leftrightarrow$ denominator height & Farey height $Q_c$, Arnold tongues & 328, 343 \\
Harmonic primes $5,7,11,13$ & entry depth $k=\mathrm{ord}_p(3)$ & 342, 343 \\
Chords & partitions of 8 in G$(6)$ & 342 \\
7-limit as hearing threshold & Euler-factor threshold $p^*\approx9.76$ & 343 \\
Resolution limit $\approx1\,\%$ & $100\xi$ as fractal step & 343 \\
\bottomrule
\end{tabular}
\end{center}

What remains open in all of these documents is why this assignment holds
at all --- whether it is analogy or identity.

% ============================================================
\section{Theorem A: There is no comma in a finite field}
% ============================================================

\begin{quote}
\textbf{Theorem A.} \BM\ Let $q=3^k$. Every element of GF$(q)^*$ has
finite order dividing $q-1$. Every chain $g, g^2, g^3,\dots$ generated by
an element $g$ returns exactly to $1$. A comma --- a remainder after an
almost-closed cycle --- is impossible in GF$(q)^*$. In $\mathbb{Q}^*$,
by contrast, no chain of fifths can close: $3^a=2^b$ has no solution
$a,b>0$.
\end{quote}

\emph{Proof.} The first part is Lagrange's theorem for the cyclic group
GF$(q)^*$ \EM. The second is unique factorisation in $\mathbb{Z}$ \EM. \qed

The content lies not in the proof but in the consequence for the corpus.
Doc.~060 places the Pythagorean comma $(3/2)^{12}/2^7$ and
$K_\text{frak}=1-100\xi$ side by side and finds the near-exact
cancellation $1.00013$. Theorem~A says \emph{where} both quantities live:
in $\mathbb{Q}^*$ and $\mathbb{R}$, not in GF$(27)^*$. The Galois structure
itself knows no remainder; all cycles close. The remainder appears only
when the finite structure is mapped onto the number line --- at the
transition from winding numbers to measured masses in SI units.

This is exactly the position of Doc.~333: $K_\text{frak}$ takes no part
in the lossless pullback, sits in $R_m$, and becomes visible only at the
SI transition. Theorem~A supplies the algebraic reason. And it is the
position of Doc.~343 Theorem~G$'$: the $1.05\,\%$ are a fractal step of
the embedding, not a property of the Galois grid.

\textbf{Corollary A$'$} \KM: The cancellation comma$\cdot K_\text{frak}\approx1$
of Doc.~060 is thus not a coincidence of two independent systems but two
measurements of the same embedding error --- once from the ratio $3/2$,
once from the ratio $4/3$ (electron coefficient, Doc.~006). The residual
$0.22$~cent is the difference of the two starting ratios, not a third
correction.

% ============================================================
\section{Theorem B: The exceptions of the 5-limit reading are Galois primes}
% ============================================================

Doc.~060 Table~2 classifies the nine Yukawa coefficients of Doc.~006:
seven are 5-limit ($2,3,5$), two are not --- strange $26/9$ carries the
prime 13, top $1/28$ carries the prime 7. Doc.~060 marks both as exceptions
without musical interpretation.

Doc.~343 Theorem~E shows independently: the prime $p$ enters the
GF$(3^k)$ hierarchy exactly at $k=\mathrm{ord}_p(3)$, and for $k\le6$
these are precisely $13\,(k=3)$, $5\,(k=4)$, $11\,(k=5)$, $7\,(k=6)$.

\begin{quote}
\textbf{Theorem B.} \BM\ All nine Yukawa coefficients of Doc.~006 have
prime factors exclusively in $\{2,3,5,7,11,13\}$. The two non-5-limit
coefficients carry exactly the primes of Galois depth $k=3$ (strange: 13)
and $k=6$ (top: 7). The prime 11 ($k=5$) occurs in no coefficient.
\end{quote}

\emph{Proof.} Direct factorisation of the nine fractions
(script B1--B5). \qed

Three remarks:

\begin{enumerate}
\item The 5-limit reading of music is the Galois reading up to depth
  $k=4$ (primes $2,3,5$; the prime 13 does enter at $k=3$, but as the
  group order $|Z_{13}|$, not as a coefficient). What appears in music as
  ``beyond the 5-limit'' is algebraically the next layer. The exceptions
  of Doc.~060 are thereby resolved.
\item Strange carries the prime 13 --- the same depth $k=3$ that carries
  the lepton orbits (GF$(27)^*\cong Z_2\times Z_{13}$, Doc.~339). Top
  carries 7 and needs $k=6$, i.e.\ GF$(729)$ --- exactly the candidate
  that R115 names for the CKM/PMNS angles.
\item Whether the quark coefficients are geometric or calibrated remains
  open after Doc.~289 \SM. Theorem~B says only: \emph{if} they are
  geometric, they lie in the right grid. That is compatibility, not proof.
\end{enumerate}

\textbf{Falsifiable consequence} \KM: A fundamental mass coefficient with
a prime factor $>13$ (or with 11) would be incompatible with Doc.~343
Theorem~E. The prediction is testable on every coefficient derived in future.

% ============================================================
\section{Observation C: Three twelves}
% ============================================================

Three quantities of order 12 occur, from three different mechanisms:

\begin{enumerate}
\item \textbf{Temperament.} The 5-limit Tonnetz is the free lattice
  $\mathbb{Z}^2$ in the exponents of 3 and 5 (modulo octave). Twelve-tone
  temperament is the quotient by the syntonic comma
  $81/80\,\widehat{=}\,(4,-1)$ and the diesis $128/125\,\widehat{=}\,(0,-3)$:
  \[
    \bigl|\mathbb{Z}^2/\langle(4,-1),(0,-3)\rangle\bigr| = |\det| = 12 .
  \]
  The 12 arises by \emph{quotient} of a torsion-free lattice \EM.
\item \textbf{Galois.} $|(\mathbb{Z}/13)^*|=12$; the element 3 has order 3
  in it (Doc.~343 Theorem~E), the four cosets of $\langle3\rangle$ are the
  four fermion sectors, their size 3 the three generations (Doc.~343
  Theorem~D, Doc.~348). Here the 12 is \emph{intrinsic torsion}: the group
  is cyclic, no quotient needed.
\item \textbf{Resolution.} The Farey denominator height
  $Q_c=\pi/\sqrt{6\varepsilon}$ (Doc.~343 Plain Language~II) gives, at
  comma resolution $\varepsilon\in\{1.36\,\%, 1.33\,\%, 1.05\,\%\}$, the
  values $Q_c\approx11.0$, $11.1$, $12.5$. Denominators above $\approx12$
  are indistinguishable at this resolution.
\end{enumerate}

\begin{quote}
\textbf{Observation C.} \SM\ The three twelves come from three different
mechanisms (quotient, torsion, resolution limit). That they agree in order
is an observation. A theorem mapping the three mechanisms onto one another
is not available.
\end{quote}

A hint that supports the observation without proving it: twelve-tone
temperament is not convention but the smallest quotient at which
$3^{12}\approx2^{19}$ closes to comma accuracy --- i.e.\ the smallest
order at which the free lattice \emph{imitates} the torsion of the finite
field. In this sense temperament is an approximation of the Galois
structure by the continuum, and Theorem~A says what is lost: exact closure.

% ============================================================
\section{Theorem D: Why the same mathematics}
% ============================================================

The question whether the assignment in the table of §1 is analogy or
identity has a precise answer in three steps, all established in the corpus:

\begin{enumerate}
\item $\tilde{T}\cdot m=1$ (Doc.~003, 010): mass is inverse period. A mass
  ratio $m_i/m_j$ is a period ratio $\tilde{T}_j/\tilde{T}_i$ --- a
  frequency ratio. This is not a metaphor but the definition of the quantity.
\item Frequency ratios on a periodic structure are characters of the
  cyclic group $S^1$; on $T^4$ they are the winding lattice $\mathbb{Z}^4$
  (Doc.~159, 189). The overtone series is the character spectrum of $S^1$;
  the mass spectrum is the character spectrum of $T^4$.
\item The $Z_3$ orbifold action and the reduction mod~3 (Doc.~336 ring
  homomorphism $\varphi(N)=N\bmod3$, Doc.~341) map the winding lattice to
  GF$(3^k)$. The Frobenius action $x\mapsto x^3$ is the sector pairing
  $k\leftrightarrow-k$ (Doc.~336 main result).
\end{enumerate}

\begin{quote}
\textbf{Theorem D.} \BM\ Harmonic analysis of a mass spectrum and Galois
classification in GF$(3^k)$ are two representations of the same object
--- the character spectrum of $T^4/Z_3$. The first reads it in
$\mathbb{Q}^*\subset\mathbb{R}$ (frequency ratios), the second in
GF$(3^k)^*$ (orbits). The passage between them is reduction mod~3 in one
direction and embedding into the continuum in the other; Theorem~A
describes what the embedding adds (the comma), Doc.~333 describes what the
reduction removes (nothing --- the pullback is lossless).
\end{quote}

\emph{Proof.} Concatenation of the three cited steps; no new proof step. \qed

The assignment in §1 is therefore not an analogy. It is the translation
table between two representations.

% ============================================================
\section{What the connections yield}
% ============================================================

The question ``what does this give us'' has five answers, each with a
concrete corpus return:

\begin{description}

\item[1. An admissibility criterion instead of curve fitting.]
  Which primes may occur at all in a fundamental ratio is fixed
  algebraically (Doc.~343 Theorem~E). Theorem~B applies this to Doc.~006
  and finds all nine coefficients in the grid. A fit would not have
  guaranteed this; the algebra forces it.

\item[2. A built-in resolution limit.]
  Harmonics has known for centuries that consonance ends at about the
  7-limit. Doc.~343 Theorem~F gives the number ($p^*\approx9.76$),
  Theorem~G$'$ the reason ($100\xi$). For mass spectroscopy this means:
  ratios with denominators $\gtrsim12$ are indistinguishable from random
  products (Doc.~342 Theorem~D). Harmonics supplies the intuition, algebra
  the limit, and both prevent over-interpretation --- the look-elsewhere error.

\item[3. A reason for the existence of commas.]
  Theorem~A: the comma is the price of embedding. $K_\text{frak}$ is not
  an extra parameter but the necessary difference between exact closure in
  GF$(27)^*$ and the number line. Doc.~060 found the agreement; Theorem~A
  says why it is no coincidence.

\item[4. A physical reading of consonance.]
  Doc.~328 shows mixing angles as coupling degrees and Arnold tongues as
  stability regions. Consonance in this reading is not aesthetics but the
  width of an Arnold tongue --- and conversely the mixing angles of the
  Standard Model are ``intervals'' with determinable consonance. This is
  the direction in which the CKM/PMNS values (R115, \SM) are to be
  approached: not as numbers but as denominator heights.

\item[5. A working intuition that holds.]
  The Euler Tonnetz was the practical precursor of $\xi$ (Doc.~060, 009).
  Theorem~D says why the precursor held: it was not an analogy but already
  the continuum representation of the object that the Galois representation
  later made exact. Whoever thinks in the Tonnetz thinks in $T^4/Z_3$ ---
  with comma.

\end{description}

% ============================================================
\section{Scope}
% ============================================================

\begin{itemize}
\item This document derives no masses and changes no numerical value in
  the corpus.
\item Observation~C is marked \SM\ and remains so until a theorem connects
  the three mechanisms.
\item The quark coefficients remain calibration values after Doc.~289;
  Theorem~B establishes compatibility, not origin.
\item The connection to auditory physiology (why the ear stops at the
  7-limit) is not a subject here; Doc.~343 Theorem~F treats the
  mathematical limit, not the biological one.
\end{itemize}

% ============================================================
\section{Summary}
% ============================================================

\begin{center}
\begin{tabular}{@{}lp{9cm}l@{}}
\toprule
Thm. & Statement & Status \\
\midrule
A & No comma in a finite field; comma = embedding error & \BM \\
A$'$ & comma$\cdot K_\text{frak}\approx1$ = two measurements of one error & \KM \\
B & All nine Yukawa coefficients in the Galois grid; 13, 7 no exceptions & \BM \\
C & Three twelves (quotient / torsion / resolution) & \SM \\
D & Harmonics and Galois: two representations of one object & \BM \\
\bottomrule
\end{tabular}
\end{center}

\section*{Verification scripts}

\texttt{2/python/Dok358\_Skripte/pruef\_358\_harmonik\_algebra.py} ---
22 assertions: A1 (GF$(3^k)^*$ explicit, $k=1..6$), A2 ($\mathbb{Q}^*$),
A3 (comma/$K_\text{frak}$), B1--B5 (prime factors of the nine coefficients),
C1--C4 (three twelves), D1--D2 (limit thresholds).

\section*{Note on the use of AI assistance}

AI assistance (Claude, Anthropic) was used for wording, LaTeX typesetting
and the creation of the verification script. The substantive claims, the
status assignment and the responsibility for the text lie with the author.
